\chapter{Conclusões}
\label{sec:conclusoes}

Este Trabalho de Conclusão de Curso dedicou-se ao desenvolvimento e à avaliação de um sistema integrado para monitoramento inteligente e controle do consumo de energia em ambientes residenciais. Buscou-se enfrentar o desafio da gestão energética em lares convencionais, muitas vezes carentes de infraestrutura de automação e de ferramentas acessíveis que proporcionem visibilidade e controle efetivo sobre o uso da energia. A solução desenvolvida demonstrou a viabilidade de conjugar um pipeline de dados robusto, orquestrado pela plataforma Dagster, com uma interface de usuário interativa e intuitiva, criada com Streamlit, visando capacitar o usuário doméstico com insights valiosos e funcionalidades de automação personalizadas. Ao longo do projeto, foram exploradas tecnologias acessíveis e métodos de análise, como o Controle Estatístico de Processo (CEP), para oferecer uma ferramenta prática e academicamente relevante.

\section{Atendimento aos Objetivos do Projeto}
\label{sec:atendimento_objetivos}

O desenvolvimento deste trabalho foi norteado por um conjunto de objetivos específicos, delineados na introdução. A seguir, detalha-se o atendimento a cada um deles:

\begin{enumerate}[a)]
    \item \textbf{Realizar uma pesquisa e análise de tecnologias existentes que permitam o monitoramento energético sem a necessidade de infraestrutura prévia:} Conforme detalhado na Seção \ref{sec:pesquisa_tecnologias}, foi conduzida uma extensa pesquisa sobre hardware de monitoramento, culminando na seleção de tomadas inteligentes Tuya Smart Wi-Fi, que operam sobre a infraestrutura Wi-Fi residencial existente, não exigindo modificações elétricas ou estruturais complexas.

    \item \textbf{Explorar soluções adequadas para residências convencionais, que não foram originalmente projetadas para automação residencial:} A escolha das tomadas Tuya e o design geral do sistema, focado em componentes \textit{plug and play} e software de código aberto, atenderam diretamente a este objetivo, validando a aplicabilidade da solução em um ambiente residencial típico sem automação prévia, conforme descrito na Seção \ref{sec:empresa} e validado na Seção \ref{sec:viabilidade_impacto}.

    \item \textbf{Desenvolver um protótipo de sistema de monitoramento e controle de fácil instalação e uso, focado na aplicação de princípios de controle estatístico de processo:} O protótipo, composto pela pipeline de dados Dagster e pela aplicação Streamlit (Seções \ref{sec:dagster_pipeline} e \ref{sec:streamlit_app}), foi desenvolvido com ênfase na facilidade de uso. A interface Streamlit, em particular, foi projetada para ser intuitiva, e a aplicação de CEP foi uma funcionalidade central da análise avançada.

    \item \textbf{Implementar algoritmos para processamento de dados, previsão de consumo e aplicação de técnicas de controle estatístico de processo:} Algoritmos para pré-processamento de dados foram implementados na pipeline Dagster (Seção \ref{subsubsec:pipeline_dagster}). Técnicas de Controle Estatístico de Processo, especificamente gráficos CUSUM, foram implementadas e detalhadas na Seção \ref{subsec:analise_cep} e disponibilizadas na interface Streamlit. A previsão de consumo, embora mencionada como objetivo inicial, foi direcionada para a seção de trabalhos futuros (Seção \ref{sec:trabalhos_futuros}), com o foco principal do desenvolvimento prático nos algoritmos de CEP.

    \item \textbf{Desenvolver uma plataforma de software intuitiva para visualização, análise dos dados e suporte ao controle estatístico de processo pelo usuário:} A aplicação Streamlit (Seção \ref{sec:streamlit_app}) cumpriu integralmente este objetivo, oferecendo múltiplas visualizações de dados, ferramentas de análise, incluindo a interface para CEP, e funcionalidades de controle de dispositivos e cenas.

    \item \textbf{Testar e validar o sistema em ambiente real, garantindo sua eficiência e usabilidade:} A validação em ambiente real, descrita na Seção \ref{sec:viabilidade_impacto} (anteriormente Seção 3.4 na Metodologia), confirmou a funcionalidade do sistema, a precisão relativa dos sensores para os propósitos do projeto, a eficácia do CEP na detecção de anomalias simuladas e a usabilidade da interface.

    \item \textbf{Promover a viabilidade de implementação em larga escala por meio da análise de custos e estratégias de adoção:} Este objetivo foi recontextualizado durante o desenvolvimento. Conforme discutido na Seção \ref{sec:viabilidade_impacto} (anteriormente Seção 3.5 na Metodologia), o foco foi direcionado para a análise de viabilidade em contexto residencial individual e para o impacto acadêmico e benefícios ao usuário, em vez de estratégias de adoção em larga escala, que extrapolavam o escopo de um TCC. Os custos de implementação para um usuário foram considerados acessíveis.

    \item \textbf{Documentar e divulgar os resultados obtidos, destacando as contribuições para a Engenharia de Controle e Automação e os benefícios sociais e ambientais proporcionados pelo sistema desenvolvido:} Esta monografia representa o principal veículo de documentação e divulgação dos resultados. As contribuições para a engenharia e os benefícios socioambientais foram discutidos ao longo do texto, especialmente nas seções de metodologia, resultados e nas considerações sobre impacto (Seção \ref{sec:viabilidade_impacto}).
\end{enumerate}

\section{Principais Contribuições}
\label{sec:contribuicoes}
As principais contribuições deste projeto incluem:
\begin{itemize}
    \item A implementação de um pipeline de dados automatizado e escalável utilizando Dagster, capaz de ingerir dados brutos da API da Tuya, processá-los com DuckDB e armazená-los em formato Parquet particionado.
    \item O desenvolvimento de uma aplicação Streamlit com uma experiência de usuário aprimorada, oferecendo visualizações claras do consumo de energia em diferentes níveis (geral da casa e por dispositivo), controle individual de dispositivos e um sistema de gerenciamento de cenas agendadas.
    \item A integração de técnicas de Controle Estatístico de Processos (CEP), como gráficos CUSUM, para detecção proativa de anomalias no consumo de energia, agregando valor à análise e permitindo ações corretivas.
    \item A orquestração do agendamento de cenas via Dagster, garantindo a execução confiável e pontual de automações configuradas pelo usuário.
    \item A demonstração da aplicação prática de Controle Estatístico de Processos (CEP) para o monitoramento de consumo energético residencial, oferecendo um método robusto para a detecção de anomalias.
    \item A validação da arquitetura proposta em um ambiente residencial convencional, confirmando sua aplicabilidade e usabilidade para usuários finais sem conhecimento técnico especializado.
\end{itemize}

\section{Limitações e Desafios}
\label{sec:limitacoes}
Durante o desenvolvimento, algumas limitações e desafios foram identificados:
\begin{itemize}
    \item A dependência da API da Tuya para a coleta de dados, que pode estar sujeita a restrições de taxa ou mudanças na política da plataforma.
    \item A necessidade de configuração manual de variáveis de ambiente para as credenciais da API, o que pode ser um ponto de atrito na implantação.
    \item A complexidade inerente à depuração de pipelines orquestrados e aplicações web em tempo real, exigindo ferramentas de monitoramento eficazes.
\end{itemize}
\section{Trabalhos Futuros}
Para expandir e aprimorar o sistema, sugere-se as seguintes direções para trabalhos futuros:
\begin{itemize}
    \item \textbf{Integração com LLMs}: Explorar a incorporação de LLMs para permitir uma interação mais natural e conversacional com a aplicação, possibilitando que os usuários consultem dados e configurem automações por meio de linguagem natural.
    \item \textbf{Comparativo com Contas de Energia da CEMIG}: Desenvolver uma funcionalidade para comparar o consumo de energia registrado pelo sistema com as contas de energia reais da CEMIG, oferecendo uma validação externa e insights sobre a precisão das medições.
    \item \textbf{Modelagem Preditiva de Consumo}: Implementar modelos de aprendizado de máquina para prever o consumo de energia futuro, auxiliando os usuários no planejamento e na otimização do uso de energia.
    \item \textbf{Gamificação e Incentivos}: Introduzir elementos de gamificação ou sistemas de incentivo para engajar os usuários na redução do consumo de energia.
    \item \textbf{Aprimoramentos de UX/UI}: Continuar refinando a experiência do usuário e a interface gráfica, incorporando feedback dos usuários e as melhores práticas de design.
    \item \textbf{Robustez do Agendamento de Cenas}: Aprimorar a lógica de tratamento de fusos horários e a robustez de cenas não recorrentes no Dagster, garantindo que as automações sejam executadas de forma ainda mais precisa e confiável em diferentes contextos.
    \item \textbf{Expansão do Monitoramento}: Integrar o sistema com outros tipos de sensores, como medidores de água ou gás, ou com medidores de energia centrais para uma visão holística do consumo doméstico.
    \item \textbf{Estudo Longitudinal}: Realizar um acompanhamento de longo prazo com um grupo de usuários para analisar o impacto real do sistema na mudança de comportamento e na economia de energia ao longo do tempo.
    \item \textbf{Alertas Proativos}: Desenvolver um sistema de notificações (e.g., via e-mail, push) para alertar os usuários sobre anomalias detectadas pelo CEP, consumo excessivo ou falhas de dispositivos.
\end{itemize}
